\section{Conclusion}
\label{sec:conclusion}

Adding a dock-velocity feedforward to a static-dock visual docking pipeline turned a reactive baseline that docks at every sway period into a method that fails at the wave periods of its operating site, and its own estimator pointed at itself. The recorded trials showed the estimator was right: the measurement was clean, the velocity state tracked the dock, and the vehicle over-swung because the feedforward was scaled for a plant gain of 1.3 to 1.6 where the command channel delivers 2.7 with a 0.4~s time constant. An oracle feedforward from the true dock velocity over-swings identically, so the limit was the open-loop law, not the estimate. Closing the feedforward as a body-velocity setpoint on the vehicle's own navigation velocity removes the over-swing at every period, docks cleanly from 20~s to 8~s with entry-plane offsets under 2~cm, and cancels a slow navigation-velocity bias that no camera-based estimator, in the world frame or the relative frame, can observe. Bounding the velocity state during marker blackout removes the metre-scale dead-reckoning that produced the apparent phantom velocity. The remaining limit at 6~s is the handoff gate, not tracking. For a deployment the order of importance is the velocity loop, then the bounded velocity state, then the choice of estimator; and since every diagnostic the system exposed pointed the wrong way, a feedforward on a moving target needs an independent measurement of the vehicle's response before its estimator is blamed.

Four extensions follow from where the results stop: size or schedule the dock filter's velocity bound with the navigation drift, or estimate the drift apart from the sway, which the high-bias miss asks for; gate the coarse handoff on the filter's relative velocity rather than a differenced pose, which the recorded 6~s trials show would latch where the present gate does not; drive the dock with a multi-sine or a measured sea state, the next test of the bounded velocity state; and carry the procedure to the vehicle, the step tests of Section~\ref{sec:design} on the real plant, the loop bandwidth set below the identified corner, and the comparison against the reactive baseline in a tank with a moving dock, where the thruster and tether lags the simulation omits will show.
